\section{Metodologia}\label{sec:metodologia}

\subsection{Fluxo Computacional Aberto e Organiza\c{c}\~{a}o da Su\'{i}te Experimental}\label{sec:met_fluxo}

Todo o desenvolvimento experimental e num\'{e}rico deste trabalho foi estruturado em um ambiente computacional integralmente aberto e reproduz\'{i}vel, eliminando a depend\^{e}ncia original de ferramentas comerciais propriet\'{a}rias---nomeadamente Synopsys HSPICE com suporte a Verilog-A (\texttt{\$table\_model}) e MATLAB---empregadas nos arquivos de refer\^{e}ncia de Wu et al.~\cite{wu2020reconfigurable}. O motor de simula\c{c}\~{a}o el\'{e}trica adotado \'{e} o \texttt{ngspice-47} (execut\'{a}vel em modo lote \texttt{ngspice\_con.exe}), acoplado a rotinas em Python (\texttt{numpy}, \texttt{pandas} e \texttt{matplotlib}) respons\'{a}veis pela gera\c{c}\~{a}o parametrizada de netlists SPICE, extra\c{c}\~{a}o automatizada de m\'{e}tricas, an\'{a}lises estat\'{i}sticas e produ\c{c}\~{a}o de figuras.

Para assegurar rastreabilidade e verifica\c{c}\~{a}o cont\'{i}nua, os experimentos foram organizados em oito m\'{o}dulos aut\^{o}nomos (\texttt{00} a \texttt{07}) sob o diret\'{o}rio \texttt{experimentos/}, executados sequencialmente pelo orquestrador \texttt{run\_all.py}, que re-simula todas as netlists e valida a integridade dos 66 artefatos gerados (bibliotecas, netlists, tabelas CSV e figuras). Complementarmente, construiu-se uma plataforma web interativa de prototipagem visual (\textit{BP-FET Studio}, em \texttt{simulador\_web/}) conectada diretamente ao \texttt{ngspice-47}. A Tabela~\ref{tab:resumo_modulos} sintetiza o escopo metodol\'{o}gico de cada componente da su\'{i}te.

\begin{table}[htbp]
\centering
\caption{Estrutura da su\'{i}te experimental em \texttt{ngspice-47} + Python e plataforma visual.}
\label{tab:resumo_modulos}
\footnotesize
\begin{tabularx}{\textwidth}{l l X l}
\toprule
\textbf{M\'{o}dulo} & \textbf{Dom\'{i}nio / An\'{a}lise} & \textbf{Escopo Metodol\'{o}gico} & \textbf{Principais Entregas} \\
\midrule
\texttt{00} & Dispositivo (DC) & Tradu\c{c}\~{a}o 4D das tabelas experimentais para SPICE nativo e valida\c{c}\~{a}o cruzada contra \texttt{.printtr0} do HSPICE. & \texttt{bp\_ari\_ngspice.lib}, 5 netlists, 2 CSVs, 4 figs. \\
\texttt{01} & Inversor 2T (DC/Tran) & Varreduras $V_{DD}$ de $1{,}0\ \text{V}$ a $0{,}2\ \text{V}$ com e sem compensa\c{c}\~{a}o de $V_{fb}$; comuta\c{c}\~{a}o transiente ($0{,}5\ \text{fF}$). & 2 netlists, 1 CSV, 2 figs. \\
\texttt{02} & Portas 4T (Tran) & Transmuta\c{c}\~{a}o din\^{a}mica em $t=20\ \text{ns}$ para \texttt{NAND}$\leftrightarrow$\texttt{NOR} ($0{,}8\ \text{V}$) e \texttt{XOR}$\leftrightarrow$\texttt{XNOR} ($1{,}5\ \text{V}$). & 2 netlists, 1 CSV, 2 figs. \\
\texttt{03} & DPA e Ru\'{i}do (DC/Estat.) & Portabilidade Python dos conjuntos DPA, Monte Carlo 3D ($N=250$) e reten\c{c}\~{a}o; varredura $5\times 5$ no \texttt{ngspice}. & 1 netlist, 3 CSVs, 4 figs. \\
\texttt{04} & SRAM 6T (Tran/DC) & Implementa\c{c}\~{a}o do Mecanismo B ($30\ \text{ns}$), curvas borboleta, SNM e \textit{Wordline Read-Assist} ($0{,}30\ \text{V}$). & 2 netlists, 1 CSV, 2 figs. \\
\texttt{05} & Full Adder 28T (Tran) & Somador polim\'{o}rfico multiest\'{a}gio com chave $(K_1, K_0)$, convers\~{a}o de n\'{i}vel e Dist\^{a}ncia de Hamming. & 1 netlist, 1 CSV, 2 figs. \\
\texttt{06} & Fronteiras (Tran escada) & Varredura de 292 pontos em tens\~{o}es n\~{a}o padronizadas dentro do dom\'{i}nio medido das LUTs. & 3 netlists, 2 CSVs, 3 figs. \\
\texttt{07} & Pareto DPA e FA (Tran) & Avalia\c{c}\~{a}o de 7 tuplas polim\'{o}rficas, Monte Carlo ($N=250$) e \textit{benchmark} de Full Adder em 3 perfis. & 4 netlists, 3 CSVs, 4 figs. \\
\texttt{Web} & Plataforma Interativa & Compilador de grafos alto/baixo n\'{i}vel para SPICE, Raio-X de transistores, oscilosc\'{o}pio DPA e testes. & Servidor HTTP, compilador, su\'{i}te E2E \\
\bottomrule
\end{tabularx}
\end{table}

\subsection{Tradu\c{c}\~{a}o do Modelo Comportamental 4D do BP SB-FET para \texttt{ngspice-47}}\label{sec:met_modelo}

No modelo original em Verilog-A (\texttt{bp\_ari.va}), o transistor ambipolar de f\'{o}sforo negro (\textit{Black Phosphorus Schottky-Barrier FET} --- BP SB-FET) \'{e} representado por um bloco \texttt{ari} de cinco terminais (\texttt{d, g, s, g1, g2}) parametrizado pela largura relativa $W$, o qual consulta duas tabelas experimentais quadridimensionais (\texttt{BP\_ARI\_\allowbreak dataN.txt}, com 1.810 linhas, e \texttt{BP\_ARI\_\allowbreak dataP.txt}, com 759 linhas) por meio da diretiva \texttt{\$table\_model}~\cite{wu2020reconfigurable}. Como o \texttt{ngspice-47} n\~{a}o executa diretamente tabelas 4D via \texttt{\$table\_model}, desenvolveu-se no M\'{o}dulo \texttt{00} (\texttt{gerar\_\allowbreak modelo\_\allowbreak e\_\allowbreak validar.py}) o subcircuito comportamental nativo \texttt{.subckt ARI d g s g1 g2 W=1.0} (\texttt{bp\_ari\_\allowbreak ngspice.lib}), em sua vers\~{a}o refinada \texttt{v2} (Superf\'{i}cie 4D).

Para suportar circuitos bidirecionais nos quais dreno e fonte trocam de papel conforme o vetor de entrada ou a chave de polaridade---como portas \textit{pass-transistor} \texttt{XOR}/\texttt{XNOR}, transistores de acesso da SRAM e seguidores de fonte no estado bloqueado---, o subcircuito define potenciais de refer\^{e}ncia bilaterais efetivos:
\begin{equation}\label{eq:bilateral_ref}
V_{g,n} = V_g - \min(V_d, V_s), \quad V_{g1,n} = V_{g1} - \min(V_d, V_s),
\end{equation}
\begin{equation}\label{eq:bilateral_ref_p}
V_{g,p} = V_g - \max(V_d, V_s), \quad V_{g1,p} = V_{g1} - \max(V_d, V_s).
\end{equation}
Essas express\~{o}es preservam exatamente as polariza\c{c}\~{o}es nominais quando $V_{ds} \ge 0$ no ramo \textit{n-FET} e $V_{ds} \le 0$ no ramo \textit{p-FET}, mantendo continuidade de corrente na invers\~{a}o de $V_{ds}$.

Sobre esse referencial, a superf\'{i}cie de transfer\^{e}ncia bidimensional $(V_{g1}, V_{g2})$ em $|V_{ds}| = 0{,}8\ \text{V}$ \'{e} reconstru\'{i}da por interpola\c{c}\~{a}o bilinear atrav\'{e}s de fun\c{c}\~{o}es-chap\'{e}u lineares por partes (\texttt{pwl}), com travamento plano externo (\textit{clamping}) em $\pm 1000\ \text{V}$ para evitar extrapola\c{c}\~{o}es n\~{a}o f\'{i}sicas e garantir converg\^{e}ncia no solucionador MNA:
\begin{itemize}
    \item \textbf{Ramo \textit{n-FET}:} $7 \times 4 = 28$ fatias experimentais cobrindo $V_{g1} \in [-0{,}1;\ 1{,}1]\ \text{V}$ (em passos de $0{,}2\ \text{V}$) e $V_{g2} \in \{0{,}2;\allowbreak\ 0{,}4;\allowbreak\ 0{,}8;\allowbreak\ 1{,}0\}\ \text{V}$, varridas em $V_{g,n} \in [-0{,}8;\ 1{,}1]\ \text{V}$.
    \item \textbf{Ramo \textit{p-FET}:} $8 \times 4 = 32$ fatias experimentais cobrindo $V_{g1} \in [-1{,}5;\ -0{,}1]\ \text{V}$ (em passos de $0{,}2\ \text{V}$) e $V_{g2} \in \{-1{,}5;\allowbreak\ -1{,}2;\allowbreak\ -1{,}1;\allowbreak\ -0{,}9\}\ \text{V}$, varridas em $V_{g,p} \in [-1{,}5;\ 0{,}0]\ \text{V}$. A inclus\~{a}o expl\'{i}cita do plano medido em $V_{g2} = -1{,}5\ \text{V}$ captura o deslocamento real de tens\~{a}o de banda plana ($\Delta V_{fb} \approx -0{,}25\ \text{V}$) presente nos dados experimentais, o qual n\~{a}o era reproduzido pela aproxima\c{c}\~{a}o escalar preliminar (\texttt{v1}).
\end{itemize}

Para modelar as regi\~{o}es de tr\'{i}odo e satura\c{c}\~{a}o em fun\c{c}\~{a}o simult\^{a}nea da porta principal e da tens\~{a}o dreno-fonte, a vers\~{a}o \texttt{v2} substitui o fator anal\'{i}tico global $\tanh(2{,}5\,V_{ds})/\tanh(2{,}0)$ da vers\~{a}o \texttt{v1} por fatores bidimensionais $f_{ds,n}(V_{g,n}, V_{ds})$ (6 fatias em $V_{g,n} \in [0{,}1;\ 1{,}1]\ \text{V}$) e $f_{ds,p}(V_{g,p}, V_{ds})$ (5 fatias em $V_{g,p} \in [-1{,}5;\ -0{,}7]\ \text{V}$), extra\'{i}dos das varreduras experimentais de sa\'{i}da, normalizados em $|V_{ds}| = 0{,}8\ \text{V}$ e estendidos como fun\c{c}\~{o}es \texttt{pwl} \'{i}mpares estritamente monot\^{o}nicas com condut\^{a}ncia \^{o}hmica finita em $V_{ds} = 0\ \text{V}$. A corrente de dreno-fonte \'{e} ent\~{a}o imposta pela fonte comportamental:
\begin{equation}\label{eq:bp_current}
I_{ds} = W \cdot \begin{cases}
I_n(V_{g1,n},\, V_{g2,s},\, V_{g,n}) \cdot f_{ds,n}(V_{g,n},\, V_{ds}), & \text{se } V(g2,s) > 0\ \text{V}, \\
I_p(V_{g1,p},\, V_{g2,s},\, V_{g,p}) \cdot f_{ds,p}(V_{g,p},\, V_{ds}), & \text{se } V(g2,s) \le 0\ \text{V},
\end{cases}
\end{equation}
acoplada em paralelo a um resistor $R_{\mathrm{gmin}} = 10^{12}\ \Omega$ para estabilidade num\'{e}rica. Importa registrar que, assim como o arquivo \texttt{bp\_ari.va} original, a biblioteca \texttt{bp\_ari\_ngspice.lib} \'{e} puramente quase-est\'{a}tica e n\~{a}o inclui capacit\^{a}ncias intr\'{i}nsecas de porta ou sobreposi\c{c}\~{a}o ($C_{gs}, C_{gd}, C_{g2s}$); por esse motivo, todas as simula\c{c}\~{o}es transientes descritas a seguir inserem capacit\^{a}ncias nodais e de carga expl\'{i}citas ($0{,}05\ \text{fF}$ a $1{,}0\ \text{fF}$).

A valida\c{c}\~{a}o do modelo foi realizada por meio de cinco netlists DC (\texttt{sim\_N\_transfer.sp}, \texttt{sim\_P\_transfer.sp}, \texttt{sim\_N\_output.sp}, \texttt{sim\_P\_output.sp} e \texttt{sim\_polarity\_gate.sp}), comparando as vers\~{o}es \texttt{v1} e \texttt{v2} diretamente com os registros oficiais \texttt{.printtr0} gerados pelo Synopsys HSPICE por meio da raiz do erro quadr\'{a}tico m\'{e}dio (RMSE), erro absoluto m\'{e}dio (MAE), erro m\'{a}ximo e coeficiente de determina\c{c}\~{a}o ($R^2$). Extra\'{i}ram-se tamb\'{e}m $I_{\mathrm{on}}$, $I_{\mathrm{off}}$, raz\~{a}o $I_{\mathrm{on}}/I_{\mathrm{off}}$ e inclina\c{c}\~{a}o m\'{i}nima de sublimiar ($SS_{\min}$) em $|V_{ds}| = 0{,}8\ \text{V}$.

\subsection{Caracteriza\c{c}\~{a}o do Inversor em Ultra-Baixa Tens\~{a}o e Portas Polim\'{o}rficas}\label{sec:met_inv_poly}

No M\'{o}dulo \texttt{01}, o inversor l\'{o}gico complementar ($W_{\mathrm{PU}} = W_{\mathrm{PD}} = 1{,}0$) foi configurado com polariza\c{c}\~{a}o nominal de banda de condu\c{c}\~{a}o (\textit{pull-down} \textit{n-FET} com $V_{g1,n} = 1{,}1\ \text{V},\ V_{g2,n} = 0{,}8\ \text{V}$ e \textit{pull-up} \textit{p-FET} com $V_{g1,p} = V_{DD} - 1{,}5\ \text{V},\ V_{g2,p} = V_{DD} - 1{,}2\ \text{V}$). Executaram-se varreduras DC em \texttt{inv\_vtc.sp} ($0 \to V_{DD}$ em passos de $2\ \text{mV}$) para $V_{DD} \in \{1{,}0;\ 0{,}8;\ 0{,}6;\ 0{,}4;\ 0{,}2\}\ \text{V}$ sob duas configura\c{c}\~{o}es:
\begin{enumerate}
    \item \textbf{Acoplamento direto:} entrada $V_{\mathrm{in}}$ aplicada sem deslocamento aos terminais de porta (\texttt{off}$_n = \text{off}_p = 0$).
    \item \textbf{Compensa\c{c}\~{a}o eletrost\'{a}tica de $V_{fb}$:} inser\c{c}\~{a}o de deslocamentos controlados $\text{off}_n = -0{,}02 - 0{,}5\,V_{DD}$ e $\text{off}_p = 0{,}5\,V_{DD} - 0{,}58$, de modo que em $V_{\mathrm{in}} = V_{DD}/2$ ambos os ramos operem alinhados no ponto de maior ganho de sublimiar ($V_{gs,n,\mathrm{eff}} = -0{,}02\ \text{V}$ e $V_{gs,p,\mathrm{eff}} = -0{,}58\ \text{V}$).
\end{enumerate}
A partir das curvas VTC extra\'{i}ram-se a excurs\~{a}o percentual $(V_{\mathrm{out}}(0)-V_{\mathrm{out}}(V_{DD}))/V_{DD}$, a tens\~{a}o de invers\~{a}o $V_M$, o ganho diferencial m\'{a}ximo $|A_v|_{\max} = \max|dV_{\mathrm{out}}/dV_{\mathrm{in}}|$ e as margens de ru\'{i}do est\'{a}ticas $NM_L$ e $NM_H$ pelos pontos de tang\^{e}ncia unit\'{a}ria ($dV_{\mathrm{out}}/dV_{\mathrm{in}} = -1$). Adicionalmente, \texttt{inv\_tran.sp} avaliou a resposta transiente ($14\ \text{ns}$) sob carga capacitiva $C_{\mathrm{load}} = 0{,}5\ \text{fF}$ para $V_{DD} = 1{,}0\ \text{V}$ e $0{,}6\ \text{V}$.

No M\'{o}dulo \texttt{02}, simularam-se em regime transiente ($40\ \text{ns}$, passo de $20\ \text{ps}$, carga $C_{\mathrm{out}} = 0{,}3\ \text{fF}$) as duas primitivas polim\'{o}rficas de 4 transistores propostas em~\cite{wu2020reconfigurable}, percorrendo exaustivamente os vetores $(A,B) \in \{00, 01, 10, 11\}$ em janelas de $5\ \text{ns}$ com comuta\c{c}\~{a}o da chave $\text{Key}$ em $t = 20\ \text{ns}$:
\begin{itemize}
    \item \textbf{\texttt{NAND} $\leftrightarrow$ \texttt{NOR} (\texttt{poly\_nand\_nor.sp}, $V_{DD} = 0{,}8\ \text{V}$):} rede superior paralela ($W=1{,}0$) conectada a \texttt{rtop}/\texttt{gtop} e rede inferior em s\'{e}rie ($W=2{,}0$, capacit\^{a}ncia no n\'{o} intermedi\'{a}rio $C_{n1}=0{,}1\ \text{fF}$) conectada a \texttt{rbot}/\texttt{gbot}. Em $\text{Key}=0$ (\texttt{NAND}), aplica-se $(r_{\mathrm{top}}, g_{\mathrm{top}}, r_{\mathrm{bot}}, g_{\mathrm{bot}}) = (0{,}8;\ -0{,}4;\ 0{,}0;\ 0{,}8)\ \text{V}$; em $\text{Key}=1$ (\texttt{NOR}), permutam-se trilhos e polaridades para $(0{,}0;\ 0{,}8;\ 0{,}8;\ -0{,}7)\ \text{V}$.
    \item \textbf{\texttt{XOR} $\leftrightarrow$ \texttt{XNOR} (\texttt{poly\_xor\_xnor.sp}, $V_{DD} = 1{,}5\ \text{V}$, $V_{\mathrm{in}} = 1{,}1\ \text{V}$):} quatro transistores id\^{e}nticos ($W=1{,}0$) em ponte sim\'{e}trica acionados por entradas diferenciais $(A,\overline{A})$ nos terminais \texttt{g} e $(B,\overline{B})$ nos terminais auxiliares \texttt{g1}, alternando $(r_{\mathrm{top}}, g_{\mathrm{top}}, r_{\mathrm{bot}}, g_{\mathrm{bot}})$ entre $(1{,}5;\ 0{,}0;\ 0{,}0;\ 0{,}8)\ \text{V}$ ($\text{Key}=0$, \texttt{XOR}) e $(0{,}0;\ 0{,}8;\ 1{,}5;\ 0{,}0)\ \text{V}$ ($\text{Key}=1$, \texttt{XNOR}).
\end{itemize}

\subsection{Seguran\c{c}a contra DPA, Robustez Estat\'{i}stica e Fronteiras de Tens\~{a}o}\label{sec:met_dpa_fronteira}

No M\'{o}dulo \texttt{03}, substitu\'{i}ram-se integralmente os \textit{scripts} MATLAB de~\cite{wu2020reconfigurable} por rotinas Python em \texttt{analisar\_dpa\_e\_robustez.py} para processar as formas de onda transientes de corrente $|I(V_{DD})|$ das fam\'{i}lias \texttt{NAND}/\texttt{NOR} e \texttt{XOR}/\texttt{XNOR} (janela de $0$ a $30\ \mu\text{s}$), os conjuntos Monte Carlo de 250 amostras de flutua\c{c}\~{a}o simult\^{a}nea nos trilhos $V_{DD}$ e $V_{SS}$ por porta (classificando os n\'{i}veis l\'{o}gicos em torno do limiar $0{,}40\ \text{V}$) e as s\'{e}ries temporais de reten\c{c}\~{a}o do n\'{o} flutuante \texttt{bblc1} at\'{e} $111\ \mu\text{s}$ ap\'{o}s pulsos de $100\ \text{ns}$, $250\ \text{ns}$ e $500\ \text{ns}$. Como valida\c{c}\~{a}o param\'{e}trica independente sobre o modelo traduzido, executou-se no \texttt{ngspice-47} a netlist \texttt{sim\_supply\_noise.sp} com uma grade $5 \times 5$ variando $V_{DD} \in [0{,}65;\ 0{,}95]\ \text{V}$ e $V_{SS} \in [-0{,}10;\ +0{,}10]\ \text{V}$ na porta \texttt{NAND} 4T, mensurando a margem diferencial est\'{a}tica $V_{\mathrm{out}}(00) - V_{\mathrm{out}}(11)$.

No M\'{o}dulo \texttt{06}, investigou-se o comportamento das portas polim\'{o}rficas sob tens\~{o}es n\~{a}o padronizadas por meio de varreduras PWL em escada com quatro c\'{e}lulas paralelas correspondentes aos vetores \texttt{00}, \texttt{01}, \texttt{10} e \texttt{11}. Com rigor metodol\'{o}gico, todas as polariza\c{c}\~{o}es foram mantidas estritamente no interior do dom\'{i}nio experimentalmente medido nas tabelas LUT: $V(g2,s)_n \in [+0{,}20;\ +1{,}00]\ \text{V}$ e $V(g2,s)_p \in [-1{,}50;\ -0{,}90]\ \text{V}$, evitando a faixa n\~{a}o medida $(-0{,}90;\ +0{,}20)\ \text{V}$. Ao todo, simularam-se 292 condi\c{c}\~{o}es operacionais distribu\'{i}das em tr\^{e}s experimentos:
\begin{enumerate}
    \item \textbf{\texttt{NAND}/\texttt{NOR} ($V_{DD} \times$ Polaridade, 70 pontos):} 7 n\'{i}veis de $V_{DD} \in \{0{,}25;\allowbreak\ 0{,}35;\allowbreak\ 0{,}50;\allowbreak\ 0{,}65;\allowbreak\ 0{,}80;\allowbreak\ 0{,}95;\allowbreak\ 1{,}10\}\ \text{V}$ combinados com 5 pares de polaridade na LUT ($\text{Sub-pol Min}\ +0{,}2/-0{,}9\ \text{V}$, $\text{Sub-pol Mod}\ +0{,}4/-1{,}1\ \text{V}$, $\text{Nom.\ NAND}\ +0{,}8/-1{,}2\ \text{V}$, $\text{Nom.\ NOR}\ +0{,}8/-1{,}5\ \text{V}$ e $\text{Sobre-pol}\ +1{,}0/-1{,}5\ \text{V}$) para $\text{Key} \in \{0, 1\}$.
    \item \textbf{\texttt{NAND}/\texttt{NOR} (Assimetria $V(g2,s)_n \times V(g2,s)_p$, 72 pontos):} grade $6 \times 6$ em $V_{DD} = 0{,}80\ \text{V}$ variando $V(g2,s)_n \in \{0{,}20;\allowbreak\ 0{,}30;\allowbreak\ 0{,}40;\allowbreak\ 0{,}60;\allowbreak\ 0{,}80;\allowbreak\ 1{,}00\}\ \text{V}$ e $V(g2,s)_p \in \{-0{,}90;\allowbreak\ -1{,}00;\allowbreak\ -1{,}10;\allowbreak\ -1{,}20;\allowbreak\ -1{,}35;\allowbreak\ -1{,}50\}\ \text{V}$ para $\text{Key} \in \{0, 1\}$.
    \item \textbf{\texttt{XOR}/\texttt{XNOR} ($V_{DD} \times V_{\mathrm{in}}$ e porta fixa, 150 pontos):} 5 n\'{i}veis de $V_{DD} \in \{0{,}60;\allowbreak\ 0{,}80;\allowbreak\ 1{,}00;\allowbreak\ 1{,}20;\allowbreak\ 1{,}50\}\ \text{V}$ $\times$ 5 amplitudes de entrada $V_{\mathrm{in}} \in \{0{,}40;\allowbreak\ 0{,}60;\allowbreak\ 0{,}80;\allowbreak\ 1{,}10;\allowbreak\ 1{,}30\}\ \text{V}$ sob polaridade nominal ($+0{,}8/-1{,}5\ \text{V}$) e sub-polarizada ($+0{,}2/-0{,}9\ \text{V}$) para ambas as chaves (100 pontos), mais 50 pontos avaliando $V_{g,\mathrm{top}} = 0\ \text{V}$ fixo em $V_{DD} \in [0{,}90;\ 1{,}50]\ \text{V}$.
\end{enumerate}
Cada vetor de sa\'{i}da normalizado $V_{\mathrm{out}}/V_{DD}$ foi discretizado em \texttt{'1'} ($\ge 0{,}75$), \texttt{'0'} ($\le 0{,}25$) ou \texttt{'X'} ($0{,}25 < V_{\mathrm{out}}/V_{DD} < 0{,}75$), classificando o ponto em \texttt{CORRETA}, \texttt{OUTRA\_FUNCAO} (n\'{i}veis digitais v\'{a}lidos realizando outra tabela booleana) ou \texttt{INDEFINIDO\_ANALOGICO} (presen\c{c}a de ao menos um estado \texttt{'X'}).

\subsection{C\'{e}lula SRAM 6T com Bloqueio de Reten\c{c}\~{a}o por Polaridade (Mecanismo B)}\label{sec:met_sram}

No M\'{o}dulo \texttt{04}, implementou-se em \texttt{ngspice-47} (\texttt{ARI\_sram\_mecB\_ngspice.sp}) uma c\'{e}lula de mem\'{o}ria SRAM 6T dotada de bloqueio de reten\c{c}\~{a}o por invers\~{a}o de polaridade (\textit{Mecanismo B}). Ao contr\'{a}rio das portas l\'{o}gicas polim\'{o}rficas que permutam os trilhos de energia, o Mecanismo B mant\'{e}m os barramentos de alimenta\c{c}\~{a}o estritamente fixos em $V_{DD} = 1{,}0\ \text{V}$ e $V_{SS} = 0{,}0\ \text{V}$, utilizando a chave criptogr\'{a}fica para inverter dinamicamente o tipo de condu\c{c}\~{a}o dos transistores dos dois inversores acoplados de forma cruzada.

Os seis transistores \texttt{ARI} seguem raz\~{o}es de largura dimensionadas para conciliar estabilidade de leitura e capacidade de escrita: \textit{pull-up} \texttt{XPU1}, \texttt{XPU2} com $W = 1{,}0$; \textit{pull-down} \texttt{XPD1}, \texttt{XPD2} com $W = 2{,}0$; e portas de acesso (\textit{pass-gates}) \texttt{XPG1}, \texttt{XPG2} com $W = 1{,}5$, com capacit\^{a}ncias nodais $C_q = C_{qb} = 1{,}0\ \text{fF}$. A Tabela~\ref{tab:sram_polarizacoes} especifica o esquema exato de polariza\c{c}\~{a}o adotado nos dois estados da chave. Em $\text{Key} = 0$ (modo autorizado), o ajuste de $V_{g2,\mathrm{PU}}$ foi sintonizado para $-0{,}2\ \text{V}$ ($V_{g2,s} = -1{,}2\ \text{V}$), evitando o deslocamento de banda plana de $V_{g2,s} = -1{,}5\ \text{V}$ identificado pelo modelo \texttt{v2}. Em $\text{Key} = 1$ (modo bloqueado), \texttt{XPU1}/\texttt{XPU2} assumem polaridade \textit{n-FET} e \texttt{XPD1}/\texttt{XPD2} assumem polaridade \textit{p-FET}, convertendo os inversores regenerativos em seguidores de fonte degradados e destruindo a biestabilidade.

\begin{table}[htbp]
\centering
\caption{Esquema de polariza\c{c}\~{a}o do Mecanismo B na c\'{e}lula SRAM 6T ($V_{DD}=1{,}0\ \text{V}$, $V_{SS}=0{,}0\ \text{V}$).}
\label{tab:sram_polarizacoes}
\footnotesize
\begin{tabularx}{\textwidth}{>{\raggedright\arraybackslash}p{2.3cm} >{\raggedright\arraybackslash}p{3.1cm} >{\raggedright\arraybackslash}p{2.9cm} >{\raggedright\arraybackslash}p{2.9cm} >{\raggedright\arraybackslash}X}
\toprule
\textbf{Estado da Chave} & \textbf{Pull-Up ($W=1{,}0$)} & \textbf{Pull-Down ($W=2{,}0$)} & \textbf{Pass-Gates ($W=1{,}5$)} & \textbf{Regime F\'{i}sico Resultante} \\
\midrule
$\text{Key}=0$ (Autorizado) & \textit{p-FET}: $V_{g1}=-0{,}5\ \text{V}$, $V_{g2}=-0{,}2\ \text{V}$ ($V_{g2,s}=-1{,}2\ \text{V}$) & \textit{n-FET}: $V_{g1}=+0{,}8\ \text{V}$, $V_{g2}=+1{,}1\ \text{V}$ & \textit{n-FET}: $V_{g1}=+1{,}8\ \text{V}$, $V_{g2}=+2{,}1\ \text{V}$ & Inversores CMOS complementares biest\'{a}veis \\
$\text{Key}=1$ (Bloqueado) & \textit{n-FET}: $V_{g1}=+1{,}8\ \text{V}$, $V_{g2}=+2{,}1\ \text{V}$ ($V_{g2,s}=+1{,}1\ \text{V}$) & \textit{p-FET}: $V_{g1}=-1{,}5\ \text{V}$, $V_{g2}=-1{,}5\ \text{V}$ & \textit{n-FET}: $V_{g1}=+1{,}8\ \text{V}$, $V_{g2}=+2{,}1\ \text{V}$ & Seguidores de fonte cruzados (colapso monost\'{a}vel) \\
\bottomrule
\end{tabularx}
\end{table}

O protocolo transiente de $30\ \text{ns}$ compreende quatro fases: (1)~escrita do bit \texttt{'1'} em $0\text{--}6\ \text{ns}$ sob $\text{Key}=0$ ($\mathrm{WL}=1{,}0\ \text{V}$ entre $2$ e $6\ \text{ns}$, $\mathrm{BL}=1{,}0\ \text{V}$, $\mathrm{BLB}=0{,}0\ \text{V}$); (2)~reten\c{c}\~{a}o ativa em $6\text{--}15\ \text{ns}$ ($\mathrm{WL}=0\ \text{V}$, crit\'{e}rio $\Delta V_{\mathrm{Key0}} = |V(q)-V(qb)| > 0{,}80\ \text{V}$ em $t=12\ \text{ns}$); (3)~ativa\c{c}\~{a}o do bloqueio em $t=15\ \text{ns}$ ($\text{Key}=1$, colapso verificado em $t=18\ \text{ns}$ com crit\'{e}rio $\Delta V_{\mathrm{Collapse1}} < 0{,}15\ \text{V}$); e (4)~tentativa de reescrita do bit oposto entre $20$ e $24\ \text{ns}$ sob $\text{Key}=1$, seguida de novo fechamento de $\mathrm{WL}$ e verifica\c{c}\~{a}o de falha de reten\c{c}\~{a}o em $t=28\ \text{ns}$ ($\Delta V_{\mathrm{Key1}} < 0{,}15\ \text{V}$).

Adicionalmente, \texttt{sram\_butterfly\_snm.sp} executa varreduras DC de meia-c\'{e}lula ($V_{\mathrm{in}} \in [0;\ 1{,}0]\ \text{V}$, passo de $5\ \text{mV}$) com \texttt{alterparam} para seis regimes combinando $\text{Key} \in \{0, 1\}$ e $\mathrm{WL} \in \{0{,}0\ \text{V}\ (\text{Hold}),\ 0{,}30\ \text{V}\ (\text{Wordline Underdrive Read-Assist}),\ 1{,}0\ \text{V}\ (\text{Read s/ Assist})\}$, com $\mathrm{BL}=1{,}0\ \text{V}$. A Margem de Ru\'{i}do Est\'{a}tica (SNM) \'{e} calculada geometricamente pelo maior quadrado inscrito no lobo da curva borboleta no referencial rotacionado de $45^\circ$:
\begin{equation}\label{eq:snm_45}
u = \frac{V_{qb} - V_q}{\sqrt{2}}, \quad v = \frac{V_{qb} + V_q}{\sqrt{2}}, \quad \mathrm{SNM} = \max_{u < 0} \frac{v_1(u) - v_2(u)}{\sqrt{2}},
\end{equation}
sobre uma grade uniforme de 2.000 pontos.

\subsection{Circuitos Combinacionais Criptogr\'{a}ficos e Mapeamento de Tuplas Polim\'{o}rficas}\label{sec:met_comb_tuplas}

No M\'{o}dulo \texttt{05}, demonstrou-se a escalabilidade multiest\'{a}gio em um Somador Completo de 1 bit (\textit{Full Adder} --- FA) de 28 transistores (\texttt{full\_adder\_polimorfico.sp}) controlado por chave de 2 bits $(K_1, K_0)$: (i)~est\'{a}gio \texttt{XXOR1} (4T em $1{,}5\ \text{V}$) controlado por $K_0$, produzindo $P = A \oplus B$ ($K_0=0$) ou $P = A \odot B$ ($K_0=1$); (ii)~bloco de condicionamento e convers\~{a}o de n\'{i}vel com quatro inversores (\texttt{2}$\times$ \texttt{BP\_INV08} a $0{,}8\ \text{V}$ e \texttt{2}$\times$ \texttt{BP\_INV11} a $1{,}1\ \text{V}$, totalizando 8T com deslocamento interno de $-0{,}27\ \text{V}$ na porta) para regenerar $(P, \overline{P})$ em ambos os dom\'{i}nios de tens\~{a}o; (iii)~est\'{a}gio \texttt{XXOR2} (4T em $1{,}5\ \text{V}$) gerando $\mathrm{Sum} = P \oplus C_{\mathrm{in}}$; e (iv)~\'{a}rvore de transporte \texttt{Cout} composta por tr\^{e}s portas polim\'{o}rficas \texttt{XG1}--\texttt{XG3} (12T em $0{,}8\ \text{V}$) controladas por $K_1$ (\texttt{NAND} em $K_1=0$ e \texttt{NOR} em $K_1=1$). Todas as quatro chaves $K_1 K_0 \in \{00, 01, 10, 11\}$ foram simuladas em $40\ \text{ns}$ sobre os 8 mintermos ($000 \to 111$, $5\ \text{ns}$ por vetor), com amostragem estabilizada em $[5i+3{,}5;\ 5i+4{,}8]\ \text{ns}$ e binariza\c{c}\~{a}o nos limiares de $0{,}75\ \text{V}$ (\texttt{Sum}) e $0{,}40\ \text{V}$ (\texttt{Cout}) para extra\c{c}\~{a}o da Dist\^{a}ncia de Hamming (HD).

No M\'{o}dulo \texttt{07}, conduziu-se um mapeamento sistem\'{a}tico do compromisso entre assimetria de corrente DPA, custo em \'{a}rea, desempenho din\^{a}mico, consumo energ\'{e}tico e robustez estat\'{i}stica em sete tuplas polim\'{o}rficas sob carga padronizada $C_{\mathrm{load}} = 0{,}3\ \text{fF}$:
\begin{itemize}
    \item \textbf{Cinco tuplas de 2 entradas (\texttt{tuplas\_\allowbreak 2entradas.sp}, ciclo de $25\ \text{ns}$ na sequ\^{e}ncia \texttt{00} $\to$ \texttt{01} $\to$ \texttt{10} $\to$ \texttt{11} $\to$ \texttt{00}):} \texttt{(NAND, NOR)} (4T, $V_{DD}=0{,}8\ \text{V}$), \texttt{(XOR, XNOR)} (4T n\'{u}cleo / 8T total com inversores de entrada, $1{,}5\ \text{V}$/$1{,}1\ \text{V}$), \texttt{(AND, OR)} (6T = est\'{a}gio \texttt{NAND/NOR} 4T seguido por inversor isolador \texttt{BP\_INV08} 2T, $0{,}8\ \text{V}$), \texttt{(NAND, XNOR)} e \texttt{(NOR, XOR)} (tuplas mistas de 6T n\'{u}cleo / 10T total em $1{,}5\ \text{V}$/$1{,}1\ \text{V}$ combinando rede \textit{pass-transistor} e ramos habilitados por chave).
    \item \textbf{Duas tuplas de 3 entradas (\texttt{tuplas\_\allowbreak 3entradas.sp}, ciclo de $45\ \text{ns}$ de \texttt{000} a \texttt{111} $\to$ \texttt{000}):} \texttt{(AOI21, OAI21)} (6T em topologia composta s\'{e}rie-paralela complementar, $0{,}8\ \text{V}$) e \texttt{(MAJ3, XOR3)} (n\'{u}cleo polim\'{o}rfico quad-sim\'{e}trico de 10T com multiplexa\c{c}\~{a}o local da fun\c{c}\~{a}o de maioria vs.\ paridade em $1{,}5\ \text{V}$, totalizando 18T ao incluir o est\'{a}gio gerador e condicionador $P = A \oplus B$).
\end{itemize}

A assimetria DPA relativa entre os dois modos de cada tupla foi padronizada conforme:
\begin{equation}\label{eq:dpa_asym}
\mathrm{DPA}_{\mathrm{asym}} = \frac{|I_{\mathrm{Modo\ 0}} - I_{\mathrm{Modo\ 1}}|}{\min(I_{\mathrm{Modo\ 0}},\, I_{\mathrm{Modo\ 1}})} \times 100\%,
\end{equation}
aplicada tanto aos picos de corrente de alimenta\c{c}\~{a}o quanto \`{a}s correntes m\'{e}dias e janelas espec\'{i}ficas de transi\c{c}\~{a}o ($0\to 1$ e $1\to 0$). Extra\'{i}ram-se igualmente a soma das larguras $\sum W$, os atrasos de propaga\c{c}\~{a}o m\'{e}dios $t_p = (t_{pLH} + t_{pHL})/2$ a $50\%$ da excurs\~{a}o, a pot\^{e}ncia m\'{e}dia $P_{\mathrm{m\acute{e}d}}$ e est\'{a}tica $P_{\mathrm{est}}$, a energia por chaveamento $E_{\mathrm{sw}}$ e o produto pot\^{e}ncia-atraso ($\mathrm{PDP} = P_{\mathrm{m\acute{e}d}} \cdot t_p$). Para avalia\c{c}\~{a}o estat\'{i}stica, \texttt{sim\_robustez\_tuplas.sp} simulou uma grade $3 \times 3$ no \texttt{ngspice-47} ($\Delta V_{DD} \in \{-10\%, 0\%, +10\%\}$ e $\Delta V_{SS} \in \{-50, 0, +50\}\ \text{mV}$), a partir da qual foram avaliadas $N = 250$ amostras Monte Carlo gaussianas ($\sigma_{V_{DD}} = 3{,}5\%$, $\sigma_{V_{SS}} = 18\ \text{mV}$).

Por fim, avaliou-se o impacto arquitetural dessas tuplas em um \textit{benchmark} de Full Adder 1-bit (\texttt{full\_adder\_perfis\_dpa.sp}, $40\ \text{ns}$, $C_{\mathrm{load}} = 0{,}3\ \text{fF}$) comparando tr\^{e}s perfis: \textbf{Perfil 1 (Baseline Comum)}, com 20T de n\'{u}cleo (\texttt{2}$\times$ \texttt{XOR} 4T + \texttt{4}$\times$ \texttt{INV} 2T + \texttt{MAJ3} 4T fixo, sem travamento l\'{o}gico); \textbf{Perfil 2 (Custo/Desempenho)}, com 28T substituindo o \textit{carry} por 3 portas \texttt{(NAND, NOR)} de 4T; e \textbf{Perfil 3 (Baixa Assimetria DPA)}, com 26T substituindo a \'{a}rvore \texttt{NAND}/\texttt{NOR} por uma \'{u}nica c\'{e}lula polim\'{o}rfica \texttt{(MAJ3, XOR3)} de 10T.

\subsection{Plataforma Interativa de Prototipagem e Simula\c{c}\~{a}o Visual (\textit{BP-FET Studio})}\label{sec:met_web}

Para democratizar a explora\c{c}\~{a}o dos circuitos reconfigur\'{a}veis sem exigir escrita manual de netlists, desenvolveu-se o ambiente web \textit{BP-FET Studio} (\texttt{simulador\_web/}), composto por um servidor HTTP/REST local em Python (\texttt{server.py} baseado em \texttt{ThreadingHTTPServer}, expondo as rotas \texttt{/api/health}, \texttt{/api/tuples}, \texttt{/api/templates} e \texttt{/api/simulate}), um compilador de grafos para SPICE (\texttt{spice\_compiler.py}), um cat\'{a}logo de oito circuitos-modelo (\texttt{templates\_catalog.py}) e uma interface gr\'{a}fica SVG/Canvas em HTML5/JavaScript.

O compilador opera em dois n\'{i}veis de abstra\c{c}\~{a}o sobre a biblioteca \texttt{bp\_ari\_ngspice.lib}:
\begin{itemize}
    \item \textbf{Alto N\'{i}vel (Portas Polim\'{o}rficas):} o usu\'{a}rio conecta blocos configur\'{a}veis entre as sete tuplas mapeadas e o inversor \texttt{INV}. O compilador insere automaticamente os est\'{a}gios de gera\c{c}\~{a}o diferencial $(A, \overline{A})$ e os inversores de convers\~{a}o de n\'{i}vel (\texttt{BP\_INV08} e \texttt{BP\_INV11}) sempre que blocos de $0{,}8\ \text{V}$ e $1{,}5\ \text{V}$ s\~{a}o encadeados em cascata mista.
    \item \textbf{Baixo N\'{i}vel (Transistores \texttt{ARI}):} permite a montagem direta de esquem\'{a}ticos el\'{e}tricos com inst\^{a}ncias de cinco terminais \texttt{(d, g, s, g1, g2)}, fontes DC e PWL dependentes ou n\~{a}o da chave $\text{Key}$, capacitores e pontas de prova.
\end{itemize}
Ao executar a simula\c{c}\~{a}o comparativa ($\text{Key}=0$ e $\text{Key}=1$), o sistema devolve \`{a} interface as s\'{e}ries reamostradas de tens\~{a}o nodal (animadas passo a passo sobre o esquem\'{a}tico), a assinatura de corrente $|I_{V_{DD}}(t)|$ para an\'{a}lise DPA, a tabela-verdade binarizada e o estado interno de cada transistor no modal \textit{Raio-X} (classificando a polaridade ativa em \textit{n-FET} ou \textit{p-FET} pelo crit\'{e}rio $V(g2, s) > 0\ \text{V}$ e o grau normalizado de condu\c{c}\~{a}o). Toda a cadeia \'{e} verificada pela su\'{i}te automatizada \texttt{test\_simulator.py}.
